\newpage
\phantomsection
\label{prf:first-derivative-test-maximum}

\begin{remark*}[Return]
\hyperref[thm:first-derivative-test-maximum]{Return to Theorem}
\end{remark*}

\begin{theorem*}[First Derivative Test: Relative Maximum]
Let $I := [a,b]$, let $f : I \to \mathbb{R}$ be continuous on $I$, and let $c \in (a,b)$ with $f$ differentiable on $(a,c)$ and $(c,b)$. If there exists $\delta > 0$ with $(c-\delta,c+\delta) \subseteq I$ and $f'(x) \geq 0$ for $x \in (c-\delta,c)$ and $f'(x) \leq 0$ for $x \in (c,c+\delta)$, then $f$ has a relative maximum at $c$.
\end{theorem*}

\begin{proof}
\textbf{Professional Standard Proof.}~\\
By the Nondecreasing theorem applied on $[c-\delta, c]$, $f'(x) \geq 0$ on $(c-\delta,c)$ implies $f$ is nondecreasing on $[c-\delta,c]$, so $f(x) \leq f(c)$ for all $x \in [c-\delta,c]$. By the Nonincreasing theorem applied on $[c,c+\delta]$, $f'(x) \leq 0$ on $(c,c+\delta)$ implies $f$ is nonincreasing on $[c,c+\delta]$, so $f(x) \leq f(c)$ for all $x \in [c,c+\delta]$. Therefore $f(x) \leq f(c)$ for all $x \in (c-\delta,c+\delta)$, establishing a relative maximum at $c$.
\end{proof}

\begin{proof}
\textbf{Detailed Learning Proof.}~\\

\textbf{Step 1.} Handle the left side of $c$. The function $f$ is continuous on $[c-\delta,c]$ and differentiable on $(c-\delta,c)$ with $f'(x) \geq 0$. By the Nondecreasing Iff Nonneg\-ative Derivative theorem, $f$ is nondecreasing on $[c-\delta,c]$. Therefore $f(x) \leq f(c)$ for all $x \in [c-\delta,c]$.

\textbf{Step 2.} Handle the right side of $c$. The function $f$ is continuous on $[c,c+\delta]$ and differentiable on $(c,c+\delta)$ with $f'(x) \leq 0$. By the Nonincreasing Iff Nonpos\-itive Derivative theorem, $f$ is nonincreasing on $[c,c+\delta]$. Therefore $f(x) \leq f(c)$ for all $x \in [c,c+\delta]$.

\textbf{Step 3.} Combine the two sides. For any $x \in (c-\delta,c+\delta)$, either $x \leq c$ (covered by Step 1) or $x \geq c$ (covered by Step 2), giving $f(x) \leq f(c)$ throughout the punctured neighbourhood.

\textbf{Step 4.} Conclude relative maximum. By definition of relative maximum, the existence of $\delta > 0$ with $f(x) \leq f(c)$ for all $x \in I \cap (c-\delta,c+\delta)$ means $f$ has a relative maximum at $c$.
\end{proof}

\begin{remark*}[Proof structure]
A direct application of the monotonicity theorems on each side of $c$ separately. The sign condition on $f'$ is converted to a monotonicity condition, which is then used to bound $f(x)$ against $f(c)$ from above on both sides.
\end{remark*}

\begin{remark*}[Dependencies]~\\
\begin{itemize}
  \item \hyperref[thm:nondecreasing-iff-nonneg-derivative]{Theorem: Nondecreasing Iff Nonneg\-ative Derivative}
  \item \hyperref[thm:nonincreasing-iff-nonpos-derivative]{Theorem: Nonincreasing Iff Nonpos\-itive Derivative}
  \item \hyperref[def:relative-maximum]{Definition: Relative Maximum}
\end{itemize}
\end{remark*}
